% Build:
%   latexmk -pdf -interaction=nonstopmode -halt-on-error -outdir=docs/.build docs/praedixa-investor-report-2026.tex
% Output PDF:
%   docs/.build/praedixa-investor-report-2026.pdf
\documentclass[11pt,a4paper]{article}

\usepackage[a4paper,margin=22mm]{geometry}
\usepackage[T1]{fontenc}
\usepackage[utf8]{inputenc}
\usepackage[french]{babel}
\usepackage{lmodern}
\usepackage{microtype}
\usepackage{hyperref}
\usepackage{xcolor}
\usepackage{booktabs}
\usepackage{tabularx}
\usepackage{longtable}
\usepackage{array}
\usepackage{ragged2e}
\usepackage{enumitem}
\usepackage{titlesec}
\usepackage{graphicx}
\usepackage{amsmath}
\usepackage{amssymb}
\usepackage{tikz}
\usetikzlibrary{arrows.meta,positioning,shapes.geometric}

\definecolor{PraedixaBlue}{HTML}{0B4F6C}
\definecolor{PraedixaTeal}{HTML}{01BA9A}
\definecolor{PraedixaGray}{HTML}{1F2933}
\definecolor{PraedixaLight}{HTML}{F3F6F8}

\hypersetup{
  colorlinks=true,
  linkcolor=PraedixaBlue,
  urlcolor=PraedixaBlue,
  pdftitle={Praedixa - Rapport Strategie Investisseurs 2026},
  pdfauthor={Praedixa}
}

\setlist[itemize]{leftmargin=*,itemsep=1.5mm,topsep=1.5mm}
\setlist[enumerate]{leftmargin=*,itemsep=1.2mm,topsep=1.2mm}
\titleformat{\section}{\Large\bfseries\color{PraedixaGray}}{\thesection}{0.6em}{}
\titleformat{\subsection}{\large\bfseries\color{PraedixaGray}}{\thesubsection}{0.6em}{}
\titleformat{\subsubsection}{\normalsize\bfseries\color{PraedixaGray}}{\thesubsubsection}{0.6em}{}

\newcommand{\pill}[1]{\colorbox{PraedixaLight}{\strut\hspace{0.45em}\textsf{\small #1}\hspace{0.45em}}}
\newcommand{\srcref}[1]{\textsuperscript{\textcolor{PraedixaBlue}{[#1]}}}
\newcommand{\exhibit}[1]{\vspace{2mm}\noindent\colorbox{PraedixaLight}{\parbox{\dimexpr\linewidth-2\fboxsep\relax}{\textbf{#1}}}\vspace{1mm}}
\newcolumntype{Y}{>{\RaggedRight\arraybackslash}X}
\setlength{\emergencystretch}{2em}

\begin{document}
\pagenumbering{roman}
\hypersetup{pageanchor=false}

\begin{titlepage}
  \centering
  {\Large \textbf{Praedixa} \par}
  \vspace{3mm}
  {\LARGE \textbf{Rapport Stratégique Investisseurs / Incubateurs 2026} \par}
  \vspace{4mm}
  {\large Format cabinet -- version de travail \par}
  \vspace{8mm}

  \pill{Ops/DAF Decision Intelligence}\hspace{2mm}
  \pill{France First}\hspace{2mm}
  \pill{Pré-revenu}\hspace{2mm}
  \pill{Pré-POC signé}

  \vfill

  \begin{tabular}{rl}
    Date: & \today \\
    Audience prioritaire: & Investisseurs Seed/Pre-Seed, incubateurs \\
    Horizon: & 18 mois \\
    Statut de référence: & Entreprise en structuration \\
    Version: & 1.0
  \end{tabular}

  \vfill

  {\small \textbf{Confidentiel} -- Document de cadrage stratégique et de due diligence opérationnelle.\par}
\end{titlepage}
\hypersetup{pageanchor=true}

\tableofcontents
\newpage
\pagenumbering{arabic}

\section{Executive Summary}
\subsection{Proposition centrale}
Praedixa se positionne comme une couche de pilotage \textbf{Ops/DAF} au-dessus des outils existants (WFM, SIRH, ERP, BI) pour résoudre un problème opérationnel récurrent: \textbf{anticiper les ruptures de couverture, choisir l'arbitrage coût/service le plus robuste, et prouver l'impact économique de chaque décision}. Le produit ne remplace pas la planification; il industrialise la qualité des décisions.

\subsection{Pourquoi le problème mérite une solution dédiée}
Trois constats soutiennent l'opportunité:
\begin{itemize}
  \item La tension de recrutement reste élevée en France (2,4 millions de projets de recrutement en 2025, dont 50,1\% jugés difficiles).\srcref{S1}
  \item Les emplois vacants demeurent significatifs (taux de 2,3\% au T3 2025 dans le privé), traduisant une difficulté structurelle d'ajustement capacité/charge.\srcref{S2}
  \item Les coûts liés aux absences et arrêts maladie restent élevés, ce qui fragilise les organisations intensives en main-d'oeuvre.\srcref{S3}
\end{itemize}

\subsection{Thèse stratégique}
La valeur de Praedixa n'est pas de \og prédire mieux \fg{} de façon abstraite, mais de \textbf{rendre les arbitrages auditables et répétables}: signal de risque à J+3/J+7/J+14, options actionnables sous contraintes, puis décision tracée et comparée au réalisé (scénarios 0\% / 100\% / réel).

\subsection{Conclusion de lecture investisseur}
Le dossier montre un positionnement défendable, une wedge claire (logistique multi-sites en France), un modèle commercial premium cohérent avec la valeur, et un plan d'exécution 18 mois réaliste pour passer de pré-revenu à premiers revenus récurrents, sous réserve d'une discipline forte sur les données, l'adoption terrain, et la gouvernance DAF/Ops.

\section{Investment Thesis}
\exhibit{Exhibit 1 -- Thèse d'investissement en cinq points}
\begin{enumerate}
  \item \textbf{Pain aigu et répétitif}: le pilotage de couverture terrain reste fragmenté, tardif, et coûteux.
  \item \textbf{Positionnement clair}: Praedixa se place en \textit{overlay décisionnel}, non en remplacement des outils transactionnels.
  \item \textbf{Time-to-value court}: démarrage via exports agrégés (CSV/XLS), sans projet SI lourd.
  \item \textbf{Monétisation lisible}: pilote payant puis abonnement annuel premium aligné sur la valeur économique captée.
  \item \textbf{Risque gérable}: risques clés identifiés (data, adoption, conformité, exécution) avec plan de mitigation opérationnel.
\end{enumerate}

\exhibit{Exhibit 2 -- Cadrage du cas de base (hypothèses explicites)}
\begin{tabularx}{\linewidth}{@{}l>{\raggedright\arraybackslash}X@{}}
\toprule
\textbf{Dimension} & \textbf{Hypothèse de travail} \\
\midrule
Stade commercial & Pré-revenu, pré-POC signé \\
Vertical wedge & Logistique et opérations multi-sites \\
Géographie prioritaire & France uniquement (18 mois) \\
Cycle d'adoption & Pilote 6--8 semaines puis conversion abonnement \\
Horizon du plan & 18 mois (avec trajectoire 36 mois en annexe) \\
\bottomrule
\end{tabularx}

\section{Contexte Marché et Signaux}
\subsection{Signaux macro pertinents}
\exhibit{Exhibit 3 -- Indicateurs de contexte France/UE}
\begin{tabularx}{\linewidth}{@{}l l >{\raggedright\arraybackslash}X@{}}
\toprule
\textbf{Indicateur} & \textbf{Valeur} & \textbf{Lecture stratégique} \\
\midrule
BMO France Travail 2025 & 2,4M projets; 50,1\% difficiles\srcref{S1} & La tension de staffing est structurelle; la couverture prévisionnelle devient critique. \\
Taux d'emplois vacants (privé) T3 2025 & 2,3\%\srcref{S2} & Le déséquilibre offre/demande alimente les arbitrages d'urgence. \\
Indemnités journalières maladie 2023 & 10,2 MdEUR\srcref{S3} & Coût macro élevé des absences, pression durable sur la capacité opérationnelle. \\
Entreprises UE avec ERP (2023) & 43,3\%\srcref{S4} & Les données sont davantage disponibles, même si hétérogènes. \\
Entreprises UE utilisant l'IA (2025) & 20\%\srcref{S5} & Le marché accepte mieux les outils analytiques décisionnels. \\
Entreprises France utilisant l'IA (2024) & 10\% (vs 6\% en 2023)\srcref{S6} & Accélération locale de l'adoption, surtout chez les grandes structures. \\
\bottomrule
\end{tabularx}

\subsection{Implication pour Praedixa}
Le signal marché n'est pas uniquement \og AI \fg{}: c'est l'intersection entre rareté de capacité, volatilité de la charge, et besoin de preuve économique pour arbitrer. Praedixa se positionne précisément sur cette intersection.

\section{Problème Économique Cible}
\subsection{Formulation du problème}
Dans une organisation multi-sites, la décision quotidienne est moins \og qui planifier \fg{} que \og comment minimiser le coût total sous contrainte de service \fg{}.

\[
\min_{a \in \mathcal{A}} \mathbb{E}[C(a)] \quad \text{s.c.} \quad \mathbb{E}[1-S(a)] \leq \varepsilon
\]

où $a$ est un scénario (HS, intérim, réallocation, ajustement de service), $C(a)$ le coût total attendu, et $S(a)$ un proxy de service.

\subsection{Conséquence business observée}
\begin{itemize}
  \item Décisions prises trop tard (J-0/J-1) donc options coûteuses.
  \item Arbitrages implicites, difficilement défendables face à la DAF.
  \item Faible capitalisation: les décisions passées ne sont pas systématiquement transformées en apprentissage.
\end{itemize}

\exhibit{Exhibit 4 -- Boucle de valeur Praedixa}
\begin{center}
\begin{tikzpicture}[
  node distance=9mm,
  box/.style={draw, rounded corners, thick, align=center, minimum width=34mm, minimum height=9mm},
  arr/.style={-Latex, thick}
]
  \node[box] (risk) {Risque\\J+3/J+7/J+14};
  \node[box, right=11mm of risk] (decide) {Scénarios\\coût/service};
  \node[box, right=11mm of decide] (act) {Action\\exécutée};
  \node[box, below=9mm of decide] (proof) {Decision log\\+ preuve};

  \draw[arr] (risk) -- (decide);
  \draw[arr] (decide) -- (act);
  \draw[arr] (act) -- (proof);
  \draw[arr] (proof) -- (risk);
\end{tikzpicture}
\end{center}

\section{Positionnement et Différenciation}
\subsection{Positionnement retenu}
\textbf{Praedixa = Ops/DAF Decision Intelligence.} Ce positionnement est choisi pour éviter le \og red ocean \fg{} des suites transactionnelles WFM, et pour ancrer la valeur sur la décision économique et l'auditabilité.

\subsection{Ce que Praedixa est / n'est pas}
\exhibit{Exhibit 5 -- Frontière produit}
\begin{tabularx}{\linewidth}{@{}>{\bfseries}l X@{}}
\toprule
Est & N'est pas \\
\midrule
Couche décisionnelle au-dessus des systèmes existants & Un nouveau logiciel de planning remplaçant WFM/SIRH \\
Moteur d'options coût/service sous contraintes explicites & Une marketplace d'intérim \\
Système de preuve (0\% / 100\% / réel) & Un dashboard passif sans protocole d'attribution \\
Démarrage via exports agrégés & Un projet SI de 6 mois avant première valeur \\
\bottomrule
\end{tabularx}

\subsection{Framework d'attributs différenciants}
\begin{itemize}
  \item \textbf{Time-to-value}: conception pour démarrer avec des exports agrégés.
  \item \textbf{Traçabilité}: paramètres de coûts et règles versionnés.
  \item \textbf{Exécution}: capacité hedging via plugin agences d'intérim (pré-réservation/exercice/annulation).
  \item \textbf{Gouvernance}: decision log exploitable en CODIR/DAF.
\end{itemize}

\section{ICP, Segmentation et Wedge France}
\subsection{ICP prioritaire (18 mois)}
\exhibit{Exhibit 6 -- ICP de référence}
\begin{tabularx}{\linewidth}{@{}lX@{}}
\toprule
\textbf{Dimension} & \textbf{Cible prioritaire} \\
\midrule
Type d'organisation & Logistique, entreposage, 3PL, opérations multi-sites intensives en main-d'oeuvre \\
Taille & 4 à 20 sites, 100 à 800 opérateurs terrain \\
Maturité SI & Outils WFM/ERP existants, exports disponibles \\
Douleur métier & HS/intérim élevés, variabilité de charge, pénalités service, arbitrages tardifs \\
Acheteurs clés & Direction Exploitation, Planning terrain, DAF \\
Champion interne & Responsable exploitation régional / PMO opérations \\
\bottomrule
\end{tabularx}

\subsection{Pourquoi France-first}
\begin{itemize}
  \item Réduction de complexité d'exécution (langue, conformité, vente, delivery).
  \item Cohérence avec le wedge initial et le contexte réseau existant.
  \item Possibilité d'industrialiser le playbook avant expansion géographique.
\end{itemize}

\subsection{Personas et message par persona}
\begin{tabularx}{\linewidth}{@{}l l X@{}}
\toprule
\textbf{Persona} & \textbf{Objectif} & \textbf{Message pivot} \\
\midrule
Dir. Exploitation & Continuité de service & \og Moins de mode pompier, plus de décisions anticipées et exécutables \fg{} \\
DAF & Justification économique & \og Arbitrages coût/service tracés, hypothèses explicites, preuve mensuelle \fg{} \\
Resp. Planning & Actionnabilité quotidienne & \og Des options classées et faisables, pas un score opaque de plus \fg{} \\
\bottomrule
\end{tabularx}

\section{TAM / SAM / SOM (France, approche assumée par hypothèses)}
\subsection{Principe}
Les estimations de marché pour le WFM/VMS varient fortement selon les périmètres. Le rapport adopte donc une méthode \textbf{bottom-up} avec hypothèses explicites et révisables plutôt qu'un TAM top-down unique.

\exhibit{Exhibit 7 -- Modèle de taille de marché (hypothèses)}
\begin{tabularx}{\linewidth}{@{}>{\RaggedRight\arraybackslash}p{3.0cm} >{\RaggedRight\arraybackslash}p{4.2cm} >{\RaggedRight\arraybackslash}p{2.0cm} Y@{}}
\toprule
\textbf{Niveau} & \textbf{Base de calcul} & \textbf{Valeur} & \textbf{Commentaire} \\
\midrule
TAM (France, cible élargie) & 3 000 organisations x 90 kEUR ACV & 270 M EUR & Hypothèse haute, inclut segments adjacents \\
SAM (wedge réaliste 18-36m) & 600 organisations x 90 kEUR ACV & 54 M EUR & Segment réellement adressable par go-to-market fondateur \\
SOM (12 mois, base) & 2 clients Pro x 90 kEUR & 180 kEUR ARR & Cohérent avec stade pré-revenu et capacité delivery \\
\bottomrule
\end{tabularx}

\subsection{Interprétation investisseur}
\begin{itemize}
  \item La taille de marché adressable est suffisante pour une trajectoire venture, mais l'exécution initiale dépend surtout de la vitesse de conversion pilote -> abonnement.
  \item Le choix d'un SOM volontairement prudent renforce la crédibilité du plan.
\end{itemize}

\section{Paysage Concurrentiel et Battlecards}
\subsection{Cartographie des alternatives}
\exhibit{Exhibit 8 -- Alternatives réelles évaluées par les clients}
\begin{tabularx}{\linewidth}{@{}l l X@{}}
\toprule
\textbf{Catégorie} & \textbf{Exemples} & \textbf{Lecture pour Praedixa} \\
\midrule
Suites WFM/HCM enterprise & UKG, Workday, SAP, Oracle & Forte couverture transactionnelle, faible focalisation sur preuve économique d'arbitrage multi-option \\
Spécialistes planning/temps & Quinyx, Skello, Horoquartz, Octime & Excellents en planification/temps, moins centrés sur attribution économique 0\%/100\%/réel \\
Solutions supply chain & Manhattan, Blue Yonder (labor management) & Fortes capacités verticales, déploiement potentiellement plus lourd \\
Statut quo & Excel + emails + routines locales & Faible coût apparent, coût réel élevé et faible auditabilité \\
\bottomrule
\end{tabularx}

\subsection{Battlecards synthétiques (6--8 acteurs)}
\exhibit{Exhibit 9 -- Grille comparative synthétique}
\begin{tabularx}{\linewidth}{@{}Y >{\centering\arraybackslash}p{2.2cm} >{\centering\arraybackslash}p{2.2cm} >{\centering\arraybackslash}p{2.2cm} >{\centering\arraybackslash}p{2.2cm}@{}}
\toprule
\textbf{Acteur} & \textbf{Planning natif} & \textbf{Moteur coût/service} & \textbf{Preuve attribuable} & \textbf{Intégration légère} \\
\midrule
Suites enterprise (UKG/Workday/SAP/Oracle) & Fort & Limité à moyen & Limité & Faible à moyen \\
Spécialistes planning (Quinyx/Skello/Horoquartz/Octime) & Fort & Limité & Limité & Moyen \\
Praedixa & Non (overlay) & Fort (coeur produit) & Fort (coeur produit) & Fort (exports agrégés) \\
\bottomrule
\end{tabularx}

\subsection{Quand Praedixa gagne / perd}
\begin{tabularx}{\linewidth}{@{}>{\bfseries}Y Y@{}}
\toprule
Gagne quand & Perd quand \\
\midrule
Le client veut prouver économiquement ses arbitrages DAF/Ops & Le client cherche d'abord à remplacer tout son WFM \\
Le besoin est de décider plus tôt, sans chantier SI lourd & Le contexte impose d'emblée des intégrations enterprise profondes \\
Le sponsor accepte discipline decision log + protocole de preuve & Le sponsor veut un \og outil magique \fg{} sans changement de rituel \\
\bottomrule
\end{tabularx}

\section{Offre, Packaging et Pricing}
\subsection{Architecture d'offre}
\exhibit{Exhibit 10 -- Packaging commercial}
\begin{tabularx}{\linewidth}{@{}l l X@{}}
\toprule
\textbf{Offre} & \textbf{Prix indicatif} & \textbf{Contenu} \\
\midrule
Pilote payant (6--8 semaines) & 35 kEUR & Base canonique, signaux J+3/J+7/J+14, scénarios v1, decision log v1, proof pack initial \\
Pro (abonnement annuel) & 90 kEUR/an & Refresh, alertes, moteur scénarios complet, proof pack mensuel, gouvernance QBR trimestrielle \\
Enterprise (abonnement annuel) & 120--150 kEUR/an & SLA renforcé, SSO, environnements, connecteurs spécifiques, gouvernance renforcée \\
\bottomrule
\end{tabularx}

\subsection{Cohérence valeur/prix}
Dans le cas de base (entreprise 10 M EUR de CA), les gains annuels potentiels raisonnables sont estimés entre 240 kEUR et 450 kEUR selon le niveau d'adoption et la qualité des données. Le pricing Pro (90 kEUR/an) laisse un ratio de valeur captée généralement favorable au client, tout en finançant une delivery premium.

\subsection{Conditions commerciales recommandées}
\begin{itemize}
  \item Contrat annuel, renouvelable.
  \item Pilote créditable à 50\% sur abonnement si signature sous 30 jours.
  \item Gouvernance de preuve incluse (rituels et reporting).
  \item Périmètre de données agrégées, pas d'individuel en phase initiale.
\end{itemize}

\section{Go-to-Market 18 Mois (France)}
\subsection{Motion commerciale retenue}
Motion hybride à dominante founder-led sales:
\begin{itemize}
  \item Outbound ciblé sur sponsors Explo/DAF.
  \item Qualification data rapide (go/no-go).
  \item Pilote packagé court.
  \item Conversion structurée en abonnement.
\end{itemize}

\subsection{Séquence standard de vente}
\exhibit{Exhibit 11 -- Séquence commerciale}
\begin{tabularx}{\linewidth}{@{}l l X@{}}
\toprule
\textbf{Étape} & \textbf{Durée cible} & \textbf{Critère de passage} \\
\midrule
Discovery call & 30 min & Douleur explicite + sponsor identifié \\
Atelier cadrage & 60--90 min & Données minimales accessibles \\
Proposition pilote & 5 jours ouvrés & Périmètre validé + protocole preuve accepté \\
Exécution pilote & 6--8 semaines & Proof pack présenté au sponsor \\
Conversion abonnement & 2--4 semaines & ROI perçu + plan de rollout validé \\
\bottomrule
\end{tabularx}

\subsection{Plan de canaux}
\begin{tabularx}{\linewidth}{@{}l X@{}}
\toprule
\textbf{Canal} & \textbf{Rôle} \\
\midrule
Outbound LinkedIn/email ciblé & Génération d'opportunités qualifiées (sponsors opérationnels et finance) \\
Partenariats écosystème & Accès à des comptes via intégrateurs/opérateurs non concurrents \\
Contenu expert & Crédibilisation de la thèse \og coût/service prouvable \fg{} \\
Références pilotes & Réduction du risque perçu à la conversion \\
\bottomrule
\end{tabularx}

\subsection{KPI GTM}
\begin{itemize}
  \item Taux de conversion discovery -> atelier.
  \item Taux atelier -> pilote signé.
  \item Taux pilote -> abonnement.
  \item Délai moyen de conversion (semaines).
  \item Ratio valeur prouvée / prix abonnement.
\end{itemize}

\section{Sales Enablement et Gouvernance de Conversion}
\subsection{Actifs commerciaux requis dès J0}
\exhibit{Exhibit 11b -- Pack d'enablement minimal (version investissable)}
\begin{tabularx}{\linewidth}{@{}l X X@{}}
\toprule
\textbf{Actif} & \textbf{Objectif} & \textbf{Niveau attendu} \\
\midrule
Deck de vente modulaire & Standardiser le discours valeur/résultats & Version CFO/Ops + version opérationnelle \\
Battlecards par concurrent & Réduire les pertes sur objections comparatives & Mise à jour mensuelle, 6--8 acteurs \\
Script de démonstration & Assurer une démo orientée décisions et non fonctionnalités & Flow unique en 30 min, variantes par persona \\
ROI calculator & Objectiver la discussion économique & Paramètres visibles, hypothèses exportables \\
Playbook objections & Réduire la variabilité commerciale founder-led & 15 objections prioritaires et réponses testées \\
\bottomrule
\end{tabularx}

\subsection{Règles opérationnelles pour maximiser le pilot-to-subscription}
\begin{enumerate}
  \item Aucun pilote sans sponsor Ops et sponsor DAF identifiés.
  \item Aucun pilote sans définition de baseline explicite.
  \item Aucune restitution pilote sans comparatif 0\% / 100\% / réel.
  \item Aucun passage en abonnement sans plan de rollout site par site.
\end{enumerate}

\subsection{Cadence de pilotage commercial recommandée}
\exhibit{Exhibit 11c -- Cadence hebdomadaire}
\begin{tabularx}{\linewidth}{@{}l l X@{}}
\toprule
\textbf{Rituel} & \textbf{Fréquence} & \textbf{Décision attendue} \\
\midrule
Pipeline review & Hebdomadaire & Priorités d'ouverture et allocation de temps fondateur \\
Deal review & Hebdomadaire & Go/no-go sur offres pilote et ajustements proposition \\
Win/loss brief & Bimensuel & Ajustements messaging, pricing, objections \\
QBR interne & Mensuel & Validation trajectoire, correction des hypothèses 12/18 mois \\
\bottomrule
\end{tabularx}

\section{Gouvernance de Preuve et Reporting Investisseur}
\subsection{Modèle de preuve de valeur}
La crédibilité de Praedixa dépend d'un protocole de preuve réplicable. Le reporting ne doit pas seulement montrer des gains; il doit montrer la chaîne causale:
\begin{enumerate}
  \item alerte générée,
  \item option retenue,
  \item action exécutée,
  \item résultat observé,
  \item comparaison contrefactuelle.
\end{enumerate}

\subsection{Format du proof pack mensuel}
\exhibit{Exhibit 11d -- Structure standard du proof pack}
\begin{tabularx}{\linewidth}{@{}l X@{}}
\toprule
\textbf{Section} & \textbf{Contenu} \\
\midrule
Résumé exécutif & Décisions clés, impact économique estimé, limites observées \\
Adoption & Taux de traitement alertes, taux d'override, temps de réaction \\
Économique & Gains bruts, coûts évités, écarts vs baseline, sensibilité des hypothèses \\
Qualité modèle & Précision top alertes, calibration, dérive éventuelle \\
Plan d'action & Ajustements règles/coûts/process, actions et owners \\
\bottomrule
\end{tabularx}

\subsection{Tableau de bord investisseur (mensuel)}
\begin{tabularx}{\linewidth}{@{}l c c c@{}}
\toprule
\textbf{KPI} & \textbf{Objectif} & \textbf{Seuil d'alerte} & \textbf{Usage comité} \\
\midrule
Pilot-to-subscription & > 50\% & < 30\% & Validation du PMF économique \\
Decision adoption rate & > 70\% & < 50\% & Vérification d'usage réel \\
Emergency avoidance & Croissant & Stagnant 2 mois & Indicateur de valeur opérationnelle \\
Time-to-first-value & < 21 jours & > 35 jours & Efficacité delivery \\
Gross margin run-rate & > 75\% & < 65\% & Scalabilité opérationnelle \\
\bottomrule
\end{tabularx}

\section{Plan Opérationnel et Roadmap 18 Mois}
\subsection{Capacité d'exécution cible}
\begin{itemize}
  \item 1 fondateur (vente + produit).
  \item 1 profil data/ML ops.
  \item 1 profil data engineering.
  \item 0,5 CSM/ops à partir des premières conversions.
\end{itemize}

\subsection{Roadmap par phase}
\exhibit{Exhibit 12 -- Roadmap 18 mois}
\begin{tabularx}{\linewidth}{@{}l X X@{}}
\toprule
\textbf{Période} & \textbf{Objectif principal} & \textbf{Livrables clés} \\
\midrule
M1--M3 & Valider le pilote de référence & Stack de preuve v1, playbook delivery, premier cas utilisable \\
M4--M6 & Convertir et standardiser & Contrat abonnement \#1, templates d'onboarding, battlecards stabilisées \\
M7--M12 & Répéter la machine commerciale & 2e et 3e pilote, 2e conversion, KPI GTM sous contrôle \\
M13--M18 & Industrialiser sélectivement & Qualité produit renforcée, support client structuré, préparation scale \\
\bottomrule
\end{tabularx}

\subsection{Gates de décision internes}
\begin{enumerate}
  \item \textbf{Gate produit}: signaux suffisamment stables pour être actionnés.
  \item \textbf{Gate commercial}: ratio pilote/abonnement soutenable.
  \item \textbf{Gate delivery}: coût de service compatible avec marge brute cible.
\end{enumerate}

\section{Scénarios Financiers et Unit Economics}
\subsection{Hypothèses structurantes}
\begin{itemize}
  \item Prix Pro: 90 kEUR/an.
  \item Pilote: 35 kEUR (partiellement crédité en cas de conversion).
  \item CAC (12 mois): 35--70 kEUR selon intensité commerciale.
  \item Marge brute SaaS hors implémentation: 80--90\%.
  \item Churn logo annuel (hypothèse prudente): 8--15\%.
\end{itemize}

\subsection{Scénarios 12 mois}
\exhibit{Exhibit 13 -- Scénarios financiers à 12 mois}
\begin{tabularx}{\linewidth}{@{}l c c c@{}}
\toprule
\textbf{Indicateur} & \textbf{Conservateur} & \textbf{Base} & \textbf{Ambitieux} \\
\midrule
Pilotes signés & 2 & 3 & 5 \\
Conversions abonnements Pro & 1 & 2 & 4 \\
ARR fin d'année & 90 kEUR & 180 kEUR & 360 kEUR \\
Revenus pilotes (bruts) & 70 kEUR & 105 kEUR & 175 kEUR \\
Priorité de gestion & Preuve & Répétabilité & Accélération \\
\bottomrule
\end{tabularx}

\subsection{Lecture unit economics}
Formule simplifiée: $\text{LTV} \approx \dfrac{\text{ACV} \times \text{marge brute}}{\text{churn}}$.

Cas illustratif: $90 \times 0{,}85 / 0{,}12 \approx 638$ kEUR de LTV théorique. Le vrai enjeu reste la capacité à maintenir la qualité de preuve et la delivery à mesure de la croissance.

\subsection{Trajectoire 36 mois (directionnelle)}
Scénario base directionnel: 25--35 clients Pro/Enterprise et 2,5--4,0 M EUR d'ARR, conditionné à:
\begin{itemize}
  \item maintien d'un churn maîtrisé,
  \item industrialisation des onboardings,
  \item conservation de la promesse de preuve économique.
\end{itemize}

\section{Risques, Due Diligence et Mitigation}
\subsection{Registre des risques prioritaires}
\exhibit{Exhibit 14 -- Registre de risques priorisés}
\begin{tabularx}{\linewidth}{@{}l c c X@{}}
\toprule
\textbf{Risque} & \textbf{Prob.} & \textbf{Impact} & \textbf{Plan de mitigation} \\
\midrule
Qualité de données insuffisante & Élevée & Élevé & DQ gates, mode dégradé, go/no-go data explicite avant engagement preuve \\
Faible adoption terrain & Moyenne & Élevé & Rituels courts imposés, alerting capé, ownership opérationnel nommé \\
Scénarios contestés par la finance & Moyenne & Élevé & Paramètres de coûts validés DAF, versioning, transparence hypothèses \\
Cycle de vente plus long que prévu & Moyenne & Moyen & Offre pilote packagée, proposition standardisée, références rapides \\
Ambiguïté résidence de données / sous-traitants & Moyenne & Élevé & Cartographie hébergement/transferts, DPA, clauses contractuelles explicites \\
Maturité juridique de la société & Moyenne & Moyen & Feuille de route juridique, documentation à jour avant montée commerciale \\
\bottomrule
\end{tabularx}

\subsection{Points de vigilance transparence}
Deux sujets doivent être traités de manière proactive dans tout processus de diligence:
\begin{itemize}
  \item \textbf{Statut de la structure}: entreprise en structuration à date de ce rapport.
  \item \textbf{Architecture d'hébergement et transferts}: nécessité de distinguer clairement site vitrine, CDN, stockage applicatif, et zones de traitement.
\end{itemize}

\subsection{Plan de mitigation 90 jours}
\begin{enumerate}
  \item Formaliser le pack conformité contractuelle (DPA, sécurité, politiques d'accès).
  \item Publier une note d'architecture de résidence des données.
  \item Standardiser le protocole de preuve pour chaque pilote.
  \item Mettre en place un reporting mensuel risques et actions.
\end{enumerate}

\section{Pourquoi Maintenant, Pourquoi Praedixa}
\subsection{Pourquoi maintenant}
\begin{itemize}
  \item Pression opérationnelle et RH toujours élevée sur les métiers terrain.
  \item Disponibilité croissante des données en entreprise, sans homogénéité parfaite.
  \item Maturité accrue des décideurs pour adopter des outils analytiques orientés impact.
\end{itemize}

\subsection{Pourquoi cette équipe/proposition}
Praedixa combine un angle produit clair, une architecture pragmatique (exports-first), et une discipline de preuve rarement présente dans les approches purement prédictives. Cette combinaison constitue le coeur de la différenciation.

\section{Conclusion}
Le dossier recommande une stratégie de construction disciplinée:
\begin{enumerate}
  \item gagner des pilotes courts à haute crédibilité,
  \item convertir sur preuve chiffrée,
  \item industrialiser sans diluer la promesse de valeur.
\end{enumerate}

La réussite dépendra moins d'une surenchère technologique que de la rigueur d'exécution sur trois axes: qualité data, adoption terrain, gouvernance économique.

\newpage
\appendix
\section{Annexe A -- Registre d'Hypothèses (détaillé)}
\exhibit{Exhibit A1 -- Hypothèses de travail utilisées dans les scénarios}
\small
\begin{longtable}{@{}p{3.3cm}p{2.1cm}p{3.7cm}p{2.1cm}p{4.2cm}@{}}
\toprule
\textbf{Hypothèse} & \textbf{Valeur} & \textbf{Justification} & \textbf{Sensibilité} & \textbf{Méthode de validation} \\
\midrule
\endfirsthead
\toprule
\textbf{Hypothèse} & \textbf{Valeur} & \textbf{Justification} & \textbf{Sensibilité} & \textbf{Méthode de validation} \\
\midrule
\endhead
Nombre de sites pilote & 4 à 6 & Taille permettant une attribution par holdout & Moyenne & Liste des sites et comparabilité pré-pilote \\
Shifts par jour & 2 & Fréquent en logistique hors 24/7 complet & Moyenne & Relevé WFM/planning \\
Heures annuelles par ETP & 1 610 h & Hypothèse prudente de capacité annuelle & Moyenne & Validation paie et RH agrégées \\
Coût chargé interne & 19,9 EUR/h & Ordre de grandeur cas de base & Élevée & Réconciliation paie/comptabilité analytique \\
Heures supplémentaires annuelles & 12 000 h & Niveau plausible dans cas de base 10 M EUR & Moyenne & Extraction paie HS \\
Majoration HS & +25\% & Règle générique hors accord spécifique\srcref{S7} & Élevée & Vérification convention/accord interne \\
Intérim annuel & 18 000 h & Hypothèse de départ pré-POC & Élevée & Facturation intérim et relevés heures \\
Part intérim urgent & 40\% & Retours terrain sur commandes tardives & Élevée & Horodatage commandes et lead time \\
Réduction last-minute & 10--20\% & Fourchette prudente pour adoption progressive & Élevée & Comparatif BAU/real piloté \\
Réduction coût non-service & 15--30\% & Dépend de la stabilisation des créneaux critiques & Élevée & KPI service et coûts associés \\
Churn annuel & 8--15\% & Fourchette SaaS B2B early stage & Élevée & Cohortes clients post-lancement \\
CAC fully loaded & 35--70 kEUR & Selon cycle et volume outbound & Élevée & Suivi CRM des coûts et temps de conversion \\
\bottomrule
\end{longtable}
\normalsize

\section{Annexe B -- Blueprint POC 6--8 semaines}
\exhibit{Exhibit B1 -- Plan d'exécution standardisé}
\begin{tabularx}{\linewidth}{@{}l X X@{}}
\toprule
\textbf{Semaine} & \textbf{Objectif} & \textbf{Livrables} \\
\midrule
S1 & Cadrage data + gouvernance & Scope pilote, dictionnaire de données, protocole décision \\
S2 & Ingestion et qualité & Pipeline initial, rapport DQ, premières corrections mapping \\
S3 & Signal de risque v1 & Alertes J+3/J+7/J+14 sur périmètre pilote \\
S4 & Scénarios coût/service v1 & Options HS/intérim/réallocation avec contraintes explicites \\
S5 & Boucle d'exécution & Mise en place decision log, rituels opérationnels \\
S6 & Attribution et preuve v1 & Comparatif 0\% / 100\% / réel sur premiers cas \\
S7 & Stabilisation & Ajustements modèles/règles/coûts avec sponsor DAF \\
S8 & Restitution et conversion & Proof pack final, décision go/no-go abonnement \\
\bottomrule
\end{tabularx}

\section{Annexe C -- Dictionnaire KPI investisseur}
\exhibit{Exhibit C1 -- KPI de pilotage de la phase initiale}
\begin{tabularx}{\linewidth}{@{}l X X@{}}
\toprule
\textbf{KPI} & \textbf{Définition} & \textbf{Usage} \\
\midrule
Pilot Win Rate & Pilotes signés / ateliers cadrage & Qualité qualification commerciale \\
Pilot-to-Subscription & Abonnements signés / pilotes terminés & Capacité à convertir sur preuve \\
Time-to-First-Value & Délai entre kick-off et première décision actionnable & Qualité delivery et adoption \\
Decision Adoption Rate & \% d'alertes avec action documentée & Réalité d'usage terrain \\
Emergency Avoidance Hours & Heures d'urgence évitées vs baseline & Mesure directe de valeur opérationnelle \\
Value-to-Price Ratio & Valeur économique attribuée / prix annuel & Défendabilité commerciale et renouvellement \\
Gross Margin excl. implémentation & (Revenu - coût infra/support run) / revenu & Scalabilité SaaS \\
\bottomrule
\end{tabularx}

\section{Annexe D -- Sources Datées}
\exhibit{Exhibit D1 -- Références utilisées}
\small
\begin{longtable}{@{}p{1cm}p{5.2cm}p{2.4cm}p{6.4cm}@{}}
\toprule
\textbf{ID} & \textbf{Source} & \textbf{Date clé} & \textbf{Lien} \\
\midrule
\endfirsthead
\toprule
\textbf{ID} & \textbf{Source} & \textbf{Date clé} & \textbf{Lien} \\
\midrule
\endhead
S1 & France Travail, Enquête BMO 2025 & 11/04/2025 & \href{https://statistiques.francetravail.org/bmo/bmopub/225146}{statistiques.francetravail.org} \\
S2 & Dares, Les emplois vacants (mise à jour) & 16/12/2025 & \href{https://dares.travail-emploi.gouv.fr/donnees/les-emplois-vacants}{dares.travail-emploi.gouv.fr} \\
S3 & DREES, Études et Résultats n\textdegree1321 (arrêts maladie) & 13/12/2024 & \href{https://www.drees.solidarites-sante.gouv.fr/publications-communique-de-presse/etudes-et-resultats/241211_ER_Arrets-Maladie}{drees.solidarites-sante.gouv.fr} \\
S4 & Eurostat, E-business integration (ERP) & 06/02/2024 & \href{https://ec.europa.eu/eurostat/documents/2995521/18404397/4-06022024-AP-EN.pdf}{ec.europa.eu/eurostat} \\
S5 & Eurostat, AI use in EU enterprises & 15/04/2025 & \href{https://ec.europa.eu/eurostat/web/products-eurostat-news/w/ddn-20250415-1}{ec.europa.eu/eurostat-news} \\
S6 & INSEE, usage de l'IA en entreprise & 2024 & \href{https://www.insee.fr/fr/statistiques/8232138}{insee.fr} \\
S7 & Ministère de l'Économie, heures supplémentaires & Consultation 2026 & \href{https://www.economie.gouv.fr/particuliers/heures-supplementaires-salaries-prive}{economie.gouv.fr} \\
S8 & Commission européenne, AI Act timeline & 2024--2026 & \href{https://digital-strategy.ec.europa.eu/en/policies/regulatory-framework-ai}{digital-strategy.ec.europa.eu} \\
S9 & Prism'emploi, Bilan 2024 et perspectives 2025 & 2025 & \href{https://www.prismemploi.eu/publication/bilan-2024-perspectives-2025/}{prismemploi.eu} \\
S10 & Document interne Praedixa, project-context.md & 2026 & \texttt{docs/project-context.md} \\
S11 & Document interne Praedixa, presentation-praedixa.md & 2026 & \texttt{docs/presentation-praedixa.md} \\
S12 & Document interne Praedixa, recherche-praedixa.md & 2026 & \texttt{docs/recherche-praedixa.md} \\
\bottomrule
\end{longtable}
\normalsize

\section{Annexe E -- Battlecards détaillées (synthèse opérationnelle)}
\exhibit{Exhibit E1 -- Battlecards express pour usage commercial}
\small
\begin{longtable}{@{}p{2.5cm}p{3.9cm}p{3.9cm}p{5.5cm}@{}}
\toprule
\textbf{Acteur} & \textbf{Forces perçues} & \textbf{Faiblesses perçues} & \textbf{Talk track Praedixa} \\
\midrule
\endfirsthead
\toprule
\textbf{Acteur} & \textbf{Forces perçues} & \textbf{Faiblesses perçues} & \textbf{Talk track Praedixa} \\
\midrule
\endhead
UKG & Forte profondeur WFM, marque installée & Déploiement et conduite du changement souvent lourds & Praedixa complète l'existant en se concentrant sur l'arbitrage économique traçable et rapide. \\
Workday & Suite enterprise intégrée et crédible finance/RH & Approche globale coûteuse pour besoin ciblé de couverture & Praedixa apporte un wedge opérationnel court terme sans replateforming SI. \\
SAP SuccessFactors & Couverture enterprise, conformité, processus structurés & Complexité d'implémentation, dépendance à l'écosystème & Positionner Praedixa comme accélérateur décisionnel connecté aux exports existants. \\
Oracle HCM & Capacités enterprise avancées et gouvernance robuste & Time-to-value initial plus long pour cas local ciblé & Proposer un pilote 6--8 semaines orienté preuve économique mesurable. \\
Quinyx & Spécialiste frontline planning/temps, UX orientée manager & Moins focalisé sur attribution financière DAF & Différencier par la couche preuve 0\%/100\%/réel et moteur coût/service. \\
Skello & Forte présence terrain en France, simplicité adoption & Positionnement prioritairement planning et organisation des équipes & Affirmer le rôle overlay: Praedixa n'entre pas en conflit, il améliore la décision économique. \\
Horoquartz / Octime & Expertise locale GTA/planning et conformité & Moins orientés simulation économique multiscénarios & Mettre en avant la capacité de ranking d'options et de journalisation décisionnelle. \\
Excel + process maison & Flexibilité apparente, coût logiciel faible & Opacité, fragilité, faible réplicabilité, surcharge managériale & Chiffrer le coût caché et démontrer la robustesse d'un protocole industrialisé. \\
\bottomrule
\end{longtable}
\normalsize

\section{Annexe F -- Checklist de due diligence (investisseur/incubateur)}
\exhibit{Exhibit F1 -- Checklist de contrôle avant comité}
\small
\begin{longtable}{@{}p{4.4cm}p{2.6cm}p{7.1cm}@{}}
\toprule
\textbf{Domaine} & \textbf{Statut cible} & \textbf{Preuve attendue} \\
\midrule
\endfirsthead
\toprule
\textbf{Domaine} & \textbf{Statut cible} & \textbf{Preuve attendue} \\
\midrule
\endhead
Statut juridique société & À jour & Kbis, statuts, délégations de signature, politique contractuelle \\
Contrats clients pilotes & Standardisés & Template MSA/POC, clauses de limitation de responsabilité, calendrier de paiement \\
Protection des données & Contrôlée & DPA, registre de traitement, politique d'accès, minimisation agrégée \\
Résidence et transferts & Clarifiés & Schéma d'architecture, localisation des données, sous-traitants et clauses associées \\
Sécurité applicative & Pilotée & Contrôles d'accès, journalisation, revue de vulnérabilités, sauvegardes \\
Gouvernance modèle & Documentée & Versioning des règles/coûts, procédure de recalibration, gestion de dérive \\
Gouvernance économique & Opérationnelle & Decision log, protocole attribution, proof packs archivés \\
Scalabilité delivery & Planifiée & Runbook onboarding, capacité équipe, SLA internes, métriques run \\
Qualité commerciale & Mesurée & Funnel CRM, win/loss, conversion pilote -> abonnement \\
Traction de référence & Crédible & Cas d'usage documentés, retours sponsors, décisions business influencées \\
\bottomrule
\end{longtable}
\normalsize

\section{Annexe G -- Plan d'action 100 jours (post-comité)}
\exhibit{Exhibit G1 -- Feuille de route exécutable immédiatement}
\begin{tabularx}{\linewidth}{@{}l X X@{}}
\toprule
\textbf{Période} & \textbf{Objectif prioritaire} & \textbf{Sorties attendues} \\
\midrule
Jours 1--30 & Fiabiliser le cadre de preuve et de conformité & Contrats pilote standardisés, matrice risques v1, protocole 0\%/100\%/réel figé, note architecture données validée \\
Jours 31--60 & Accélérer la machine commerciale fondatrice & Pipeline priorisé, deck CFO/Ops finalisé, battlecards opérationnelles, 2--3 ateliers cadrage qualifiés \\
Jours 61--100 & Convertir la traction en revenus récurrents & 1er pilote structuré, preuve intermédiaire partagée, plan de conversion abonnement, roadmap produit Q+1 validée \\
\bottomrule
\end{tabularx}

\vspace{2mm}
\noindent\textbf{Principe directeur}: chaque jalon doit produire un actif réutilisable (template, process, métrique) afin d'augmenter la vitesse d'exécution sans dégrader la qualité de preuve.

\vspace{4mm}
\noindent\textbf{Contact} -- Steven Poivre -- \href{mailto:steven.poivre@outlook.com}{steven.poivre@outlook.com}

\end{document}
